\section{Introduction}
\label{sec:introduction}

\subsection{Background}
Accurate demand forecasting underlies nearly every operational decision a retail organization makes: inventory replenishment, staffing, pricing, promotional planning, and supply-chain coordination all depend on a credible estimate of how much of a given product will sell over a future horizon~\cite{fildes2022retail}. Retail sales series are shaped simultaneously by long-run trend, weekly and annual seasonality, promotional events, price changes, and idiosyncratic noise, and different products within the same organization can exhibit qualitatively different demand patterns, from smooth, high-volume staples to sparse, intermittent demand for slow-moving or newly introduced items~\cite{croston1972}. No single forecasting technique performs best across this diversity of patterns, a conclusion repeatedly reinforced by large-scale forecasting competitions such as the M4 and M5 competitions~\cite{makridakis2020m4,makridakis2022m5}.

\subsection{Business Motivation}
For small and mid-sized retailers in particular, the practical obstacle is rarely the absence of a suitable algorithm in the academic literature; it is the absence of a system that (i) works across heterogeneous, often messy uploaded datasets without bespoke data-engineering effort, (ii) selects an appropriate model automatically rather than requiring in-house data-science expertise, and (iii) produces an output that a non-technical manager can trust and act upon. A forecast that is numerically accurate but delivered as an opaque number, with no explanation of why a particular model was chosen or how confident the system is, is of limited practical value to a retail operations team.

\subsection{Challenges in Retail Forecasting}
Three challenges recur throughout this problem space. First, \emph{heterogeneity}: uploaded datasets vary in file format, column naming convention, date format, sampling frequency, and completeness, and a forecasting pipeline intended for general use must detect and normalize all of these automatically. Second, \emph{model selection under uncertainty}: because no model dominates across all series, a credible system must train a diverse set of candidates and adjudicate among them using an evaluation procedure that is robust to a single badly-fit outlier model distorting the comparison. Third, \emph{intermittent and small-sample regimes}: many retail items sell too infrequently, or a dataset is too short, for computationally expensive models such as deep neural networks to be trained reliably, and a well-designed system should recognize this rather than train models likely to overfit or fail silently.

\subsection{Machine Learning and Deep Learning Approaches}
Statistical time-series methods such as ARIMA and Holt-Winters exponential smoothing~\cite{boxjenkins1970,winters1960holtwinters} remain strong baselines, particularly on short series, and modern additive-regression tools such as Prophet~\cite{taylor2018prophet} extend this family with automatic seasonality and holiday-effect modeling. Tree-based machine learning models, most notably random forests~\cite{breiman2001randomforest} and gradient-boosted trees~\cite{chen2016xgboost}, have proven competitive on tabular, feature-engineered representations of time series, while recurrent architectures such as LSTM and GRU networks~\cite{hochreiter1997lstm,cho2014gru}, convolutional sequence models such as temporal convolutional networks~\cite{bai2018tcn}, and attention-based transformer architectures~\cite{vaswani2017attention} have extended the state of the art on longer, richer series, at the cost of requiring substantially more training data to avoid overfitting~\cite{lim2021survey}.

\subsection{The Need for Explainable AI}
A parallel and increasingly important line of work addresses the interpretability of forecasting and machine learning systems more generally. Post-hoc explanation techniques such as SHAP~\cite{lundberg2017shap} and LIME~\cite{ribeiro2016lime} have become common tools for exposing feature-level contributions to individual predictions, and the broader emergence of instruction-following large language models~\cite{brown2020gpt3,devlin2019bert} has opened a complementary avenue: rather than only exposing numerical attributions, a system can generate a natural-language explanation of its own output. When such a language model is grounded in retrieved domain knowledge and in the system's own deterministic computations, using a retrieval-augmented generation architecture~\cite{lewis2020rag}, the risk of the explanation contradicting or fabricating facts about the underlying forecast can be substantially reduced. This is the approach adopted in the system described in this paper.

\subsection{Research Motivation}
The motivation for this work is therefore not to propose a new forecasting algorithm, but to design and implement a complete, end-to-end system that (i) is dataset-agnostic across retail, transactional, and general time-series data; (ii) trains and fairly adjudicates a broad ensemble of nineteen models spanning statistical, machine learning, deep learning, and intermittent-demand families; (iii) degrades gracefully rather than failing when a dataset is too small or a candidate model is unsuitable; and (iv) closes the interpretability gap using a combination of native feature-importance scores and a grounded, retrieval-augmented language model layer, rather than treating the forecast as an opaque endpoint.

\subsection{Contributions}
The contributions of this paper are summarized as follows.
\begin{itemize}
    \item A unified orchestration layer that trains nineteen forecasting models (seven statistical, four classical machine learning, seven deep learning, and one intermittent-demand model) concurrently on an arbitrary uploaded time series, using an adaptive walk-forward cross-validation scheme whose fold count scales with the available data.
    \item A robust, multi-criteria model-selection score that normalizes root-mean-squared error, mean absolute error, mean absolute percentage error, forecast stability, computation time, and coefficient of determination using a Tukey-fence-based robust scale, preventing a single catastrophically mis-fit model from distorting the ranking of the remaining candidates.
    \item A frequency-aware gating mechanism that withholds deep learning models from datasets below an empirically justified minimum row count, communicating this decision explicitly to the end user rather than training and reporting on an unreliable model.
    \item Integration of real exogenous features (price and promotional indicators), when present in the uploaded dataset, into the gradient boosting, random forest, and seasonal autoregressive models, with an explicit and conservative policy for extending these features into the unobserved forecast horizon.
    \item A retrieval-augmented, large-language-model explanation layer, grounded simultaneously in a curated retail-domain knowledge base and in the pipeline's own deterministic metrics, that provides a conversational question-answering interface and a structured business insights report, with explicit prompt-level constraints preventing the language model from contradicting the system's own computed labels.
    \item An empirical evaluation of the complete pipeline on a real retail transaction dataset, reporting accuracy, computation time, and the practical divergence between single-metric and multi-criteria model selection.
\end{itemize}

\subsection{Paper Organization}
The remainder of this paper is organized as follows. Section~\ref{sec:literature} reviews related work in retail and time-series forecasting, machine learning and deep learning approaches, and explainable artificial intelligence. Section~\ref{sec:problem} states the problem addressed and the corresponding research gap. Section~\ref{sec:proposed} describes the proposed system architecture and workflow. Section~\ref{sec:dataset} describes the dataset used for evaluation and the automated preprocessing pipeline. Section~\ref{sec:methodology} details every implemented forecasting model, including governing equations. Section~\ref{sec:xai} describes the explainability layer. Section~\ref{sec:architecture} presents the system architecture diagrams. Section~\ref{sec:experimental} describes the experimental setup. Section~\ref{sec:results} reports quantitative results. Section~\ref{sec:discussion} interprets these results. Section~\ref{sec:future} outlines future work, and Section~\ref{sec:conclusion} concludes the paper.
